\documentclass{article}
\usepackage[left=1cm,right=1cm,top=1.5cm,bottom=1.5cm,includeheadfoot]{geometry}
%%\usepackage[utf8]{inputenc}
\usepackage[ngerman]{babel}
\usepackage[table]{xcolor}
\usepackage{nofloat}
\usepackage{listings}
\usepackage{fancyhdr}
\usepackage{pict2e}
\usepackage{varioref}
\usepackage{graphicx}
\usepackage{pdfpages}
\usepackage{luacode}
\usepackage{subcaption}
\usepackage[pdfborderstyle={/S/U/W 1},colorlinks=true, linkcolor=blue]{hyperref}
\usepackage{cleveref}
\usepackage{wasysym}
\usepackage{longtable}
\usepackage{tikz}
\usepackage{circuitikz}

\newcommand{\myfig}[3]{%
    \begin{nofloat}{figure}
        \centering
        \includegraphics[width=0.5\linewidth]{#1}
        \caption{#2}
        \label{#3}
    \end{nofloat}
}

\hypersetup{
 pdftitle={Bahnsteuerung Inbetriebnahme},
 pdfauthor={Roman Pollak},
 pdfsubject={Bahnsteuerung Inbetriebnahme}
}

\pagestyle{fancyplain}
\fancyhf{}
\lhead{\vspace{.1cm}{\includegraphics[scale=.2]{Pollak_eng}}}


\setlength{\parindent}{0pt}

\begin{document}



\begin{luacode*}
function dot_write(text)
tmp_name=os.tmpname()
--tex.print(" Roman 123"..tmp_name)
f=io.open(tmp_name..".dot",'w')
f:write(text)
f:flush()
f:close()
os.execute("dot "..tmp_name..".dot -Tpdf -o"..tmp_name..".pdf")
--#tex.print("\\newline{}\\includegraphics[width=.5\\textwidth]{"..tmp_name..".pdf}")
tex.print("\\includegraphics[width=.5\\textwidth]{"..tmp_name..".pdf}")
end
\end{luacode*}




\title{\Huge{Bahnsteuerung Inbetriebnahme / Anleitung}}
\author{Roman Pollak} 
\date{\today}
\maketitle
\rhead{\thepage}

\tableofcontents

\newpage

\section{Inbetriebnahme der Bahn}

\begin{nofloat}{figure}
        \centering
\begin{luacode*}
  dot_write([[ digraph G {

    subgraph cluster_1 {
    node [style=filled,shape=box];
    b0 -> b01 -> b1  -> b2 -> b3 -> b4 -> b5 -> b6 -> b7 -> b8 -> b9 -> b10 -> b11;
#    b3 -> b4 [label="  Messung Gut"];
#    b3 -> b5 [label="  Messung Nicht Gut"];
#    b5 -> b3;
#    b4 -> b44;
#    b44 -> b6;
#    color=blue
  }
  
  b0 [label="Bahn mechanisch Aufbauen"]
  b01 [label="Tragseile spannen\nBeachten das keine Kurzschlüsse entstehen\n Kabinen noch nicht aufhängen!"]
  b1 [label="Kabel an die Steuerung anschliessen (wenn nicht angeschlossen)"]
  b2 [label="Stromversorgung der Kabinen (Seil) ausstecken! "]
  b3 [label="Steuerung einschalten und die Debug Taste drucken cca 30 sekunden.\nUSB Dongel erst nach dem Einschalten in PC einstecken!!"]
  b4 [label="mit 'readport' überprüfen der Endschalter und Tasten"]
  b5 [label="mit 'drive 50' danach mit 'checkpos' überprüfen, ob Lichschranke richtig funktioniert\nund die Schrite gezählt werden"]
  b6 [label="Steuerung ausschalten,Motorkabel der Steuerung ausstecken,USB Dongle ausstecken\nMotoren mit Netzteil verbinden."]
  b7 [label="Kabinen einhängen\nüberprüfen mit Netzteil ob Kabinen gut fahren."]
  b8 [label="Stromversorgung der Kabinen (Seil) anschliessen\nMotoren an Steuerung anschliessen\nden Fahrtschalter auf 'Hier' stellen."]
  b9 [label="Nach dem einschalten der Steuerung sollte ein Test alles Servo von der Steuerung durchgeführt werden.\nWenn nicht, evtl ist USB Dongle noch mit PC angeschlossen"]
  b10 [label="Von der Tallstation aus sollte die Linke Kabine langsam nach unten fahren bis zum Endschalter.\nWenn die Rechte Kabine nach unten fährt, überprüfen ob der Fahrtschalter auf 'Hier' steht"]
  b11 [label="Start Taste betätigen und Gute Fahrt!!"]
  
  
  
}
]])

\end{luacode*}
\caption{Ablauf der Installation, folgend Details zu der Installation}
        \label{fig:mess}
\end{nofloat}






%%\newpage

\section{Anschlüsse der Steuerung}

\myfig{controller_new}{Anschlüsse der neuen Steuerung}{fig:ctrl}
\myfig{seilbahn}{Anschlüsse der alten Steuerung}{}
\myfig{subcontroller}{Anschlüsse an der Kabinensteuerung}{}
 
\section{ID's der Komponenten}
\myfig{Bahn_def}{ID's der Kabinen-/Towersteuerung}{fig:id}

\begin{table}[h]
    \centering
    \begin{tabular}{|c|c|c|}
        \hline
        \textbf{Lokation}&\textbf{Servo ID} & \textbf{Servo ID im C Code}\\
        \hline
        Innenseite & 1 & 0\\
        Tower Rechts & &\\
        \hline
        Aussenseite & 2 & 1\\
        Tower Links & &\\
        \hline
       
    \end{tabular}
\end{table}

\section{Debug Shell}
Durch drucken der ``Debug Shell'' Taste \autoref{fig:ctrl}, gelangt man in debug shell.

\begin{verbatim}
Debug shell                                                                                         
Be very carefull what you are doing!! Beat!! Lucas you too!!                                        
                                                                                                    
>             
\end{verbatim}

Mit ``help'' zeigt die shell alle möglichen Befehle.
\begin{verbatim}
>help
drive 
set 
readvar 
readport 
checkpos 
s2servo 
drive2 
home 
testdoor 
help 
? 
\end{verbatim}

Mit ``readvar'' werden die konfigurierbare Parameter angezeigt und deren im eeprom gespeicherten Wert. 

\begin{verbatim}
>readvar
SPEED MIN[minspeed] 55 MAX[maxspeed] 255
Sending Commands[cmdcnt] 2 times
Tower Stop[maststop] at 1465 position
Runs per button[fahrtcnt] 4
Total Button Press 16
Total runs 60
\end{verbatim}
Die Angaben in den eckigen Klammern, sind die Name für den ``set'' Befehl! Im folgenden Beispiel, wird die Anzahl der Fahrten per Knopfdruck auf 2 gesetzt. Diese Wert wird gleich im eeprom gespeichert und von der Software gleich verwendet (kein reboot notwendig).

\begin{verbatim}
>set fahrtcnt 2
SPEED MIN[minspeed] 55 MAX[maxspeed] 255
Sending Commands[cmdcnt] 2 times
Tower Stop[maststop] at 1465 position
Runs per button[fahrtcnt] 2
Total Button Press 16
Total runs 60
\end{verbatim}

Mit ``readport'' kann man die Funktionsweise der Endschalter, der ``Go'' Taste, überprüfen. Hier ist zu beachten, dass der Wert beim betätigen der entsprechenden Schalter, ändern soll!

\begin{verbatim}
>readport
Portc is 3d
Portc is 3d
Portc is 3d
..
Portc is 3d
Portc is 3d
Portc is 3c
Portc is 3c
..
Portc is 3c
Portc is 3c
Portc is 3d
Portc is 3d
..
Portc is 3d
Portc is 3d
Portc is 39
Portc is 39
\end{verbatim}

Durch drucken irgendeine Taste (der so gennante ``Any Key'') wird die Ausgabe wieder gestoppt.



Mit ``s2servo'' ist es möglich im Debug shell servos einzeln anzusteuern. zb mit ``s2servo 1 0 open'' sollte auf der rechten Kabine die Innentüre aufgehen. Mit ``s2servo 1 0 close'' wieder zu gehen.




\section{Anhang}
\subsection{Funktionsweise der Übertragung}
\myfig{Comm}{Seriale Übertragung zu den Kabinen und Tower}{}
\myfig{Comm_princip}{Funktionsweise der Detektierung}{}

\subsection{Programmierung der Kontroller mit Avrdude}
Die Kontroller kann man mit dem avrdude programmieren resp flashen. Flashen der Haupsteuerung erfolgt mit folgenden Befehl:
\begin{verbatim}
#avrdude -c usbtiny -p m168 -U flash:w:seil.hex -U eeprom:w:seil.eep
\end{verbatim}
Dabei werden alle gespeicherte Werte leider überschrieben! \\[1cm]
Programmierung der subkontroller ist etwas komplexer, da jeder subkontroller seine eigene ID besitzt. Diese ist im programm konfiguriert.
Das entsprechende ``Makefile'' des subkontrollers, erzeugt die hex dateien für alle drei ID. Die entsprechenden ID's kann man der \autoref{fig:id} entnehmen.
\begin{verbatim}
#avrdude -c usbtiny -p t841 -U flash:w:main_2.hex
\end{verbatim}
Im oberen Beispiel ist wird der Subkontroller für mit der ID 2 geflashed. ``t841'' ist die Bezeichnung des ATtiny841 im avrdude.


\section{Anmerkung}
Gefunde Fehler könnt Ihr behalten.


\end{document}
